\documentclass[]{article}
\usepackage[margin = .75in]{geometry}

\begin{document}
\section{Larger documents}
- Keep things small

- Keeping the preamble clean and organized as it grows
\section{Collaboration}
\section{ToC, Lo$\ast$}
\section{Bibliography \& Citations}
- BibLaTeX vs. BibTeX — why BibLaTeX is usually preferred now

- Citation styles (numeric, author-year, Chicago, etc.)

- Managing `.bib` files with tools like Zotero or JabRef

- Multiple bibliographies, annotated bibliographies

- Ctrl + Space within \verb|\cite|
\section{Organizing your project}
\verb|include|, \verb|input|
\section{Modular LaTeX}
\subsection{Subfiles}
\subsection{Standalone document class}
\subsubsection{Complicated TikZ, TikZ-cd, and Tables}
A good use case for the standalone document class is to use the \verb|crop| option. If you have a complicated TikZ picture or table, this can add a lot of compilation time to your document. Instead of including the raw text, we can instead create a new project with \verb|\documentclass[crop = true]{standalone}| to compile a minimally sized pdf that we can save and import instead. See \verb|/Examples/standalone_SmallExample.tex| as an example.
\section{Index \& Glossaries}
- Generating an index with `makeidx` or `imakeidx`

- Glossaries and acronym management with the `glossaries` package
\section{Beamer projects}
The author of this guide has some \textit{opinions} about how you should set up a sinlge Beamer project for all presentations someone will do during their graduate experience. It will be detailed in a separate project. Also the Beamer class has some cool unique commends that we can use and would clutter a general LaTeX guide.
\end{document}